\chapter{函数}\label{ch:g10:functions}

函数是一台机器，接收一个数，输出一个数，
对同一输入总是给出相同的输出。函数是你今后几年要学习的
分析学的核心对象
（自\cref{ch:g11:quad}起）；本章建立基本词汇 ---
定义域、像、图像、变化情况 --- 并训练
从图像中读取信息的基本技能。

\section{词汇：像与原像}

\begin{definition}[函数]\label{def:g10:functions:function}
设 $D$ 为实数集，\emph{函数}\index{函数} $f$ 为定义在
$D$ 上的对应规则，对每个 $x \in D$，恰好对应一个实数，
记为 $f(x)$，称为 $x$ 的\emph{像}\index{像}。
集合 $D$ 称为\emph{定义域}\index{定义域}，即 $f$ 的定义域。记为
\[
f : x \longmapsto f(x).
\]
若 $f(x) = y$，则 $x$ 是\emph{原像}\index{原像}，其像为 $y$。
\end{definition}

\begin{remark}
定义域中每个 $x$ 恰有一个像，但一个数 $y$ 可以有
若干个原像，也可以没有原像。对 $f(x) = x^2$，$3$ 的像是 $9$，而
$9$ 有两个原像，即 $3$ 与 $-3$；数 $-4$ 没有原像。
\end{remark}

\begin{example}[求定义域]\label{ex:g10:functions:domain}
当函数由一个式子给出时，其定义域是使该式有意义的全体 $x$ 构成的
集合。
\begin{itemize}
  \item $f(x) = 3x^2 - 5x + 1$ 对每个 $x$ 都有意义，定义域为
        $\R$.
  \item $g(x) = \dfrac{1}{x - 2}$ 要求 $x - 2 \neq 0$，定义域为
        除 $2$ 外的全体实数。
  \item $h(x) = \sqrt{x - 3}$ 要求 $x - 3 \geq 0$，定义域为
        $\intco{3}{+\infty}$.
\end{itemize}
\end{example}

\begin{example}[计算像与原像]\label{ex:g10:functions:images}
设 $f(x) = x^2 - 4x + 3$，定义在 $\R$ 上。

\emph{$5$ 的像：}代入 $x = 5$：
$f(5) = 25 - 20 + 3 = 8$.

\emph{$3$ 的原像：}解方程 $f(x) = 3$：
\[
x^2 - 4x + 3 = 3
\quad\Longleftrightarrow\quad
x^2 - 4x = 0
\quad\Longleftrightarrow\quad
x(x - 4) = 0,
\]
故 $3$ 的原像为 $0$ 与 $4$。
\end{example}

\section{函数的图像}

\begin{definition}[图像]\label{def:g10:functions:graph}
在坐标系中，\emph{图像}\index{图像}（或\emph{曲线}）是
$f$ 的全体点 $(x, f(x))$（$x$ 取遍定义域）组成的集合。
换言之，点 $(x, y)$ 在图像上，当且仅当 $y = f(x)$。
\end{definition}

\begin{omfigure}
\begin{tikzpicture}
\begin{axis}[omaxis,
  width=8.6cm, height=6.0cm,
  xmin=-2.6, xmax=4.6, ymin=-3.4, ymax=4.4,
  xtick={-2,-1,1,2,3,4}, ytick={-3,-2,-1,1,2,3,4}]
  \addplot[omDef, thick, domain=-2.2:4.2, samples=120]
    {0.25*x^3 - 0.75*x^2 - x + 2};
  \draw[dashed, omThm] (axis cs:3,0) -- (axis cs:3,-1)
    -- (axis cs:0,-1);
  \fill[omThm] (axis cs:3,-1) circle (1.5pt);
  \node[omThm, font=\small, below right] at (axis cs:3,-1.05)
    {$(3,\,f(3))$};
\end{axis}
\end{tikzpicture}

{\small 在图像上读取像：从横轴上 $x = 3$ 出发，
竖直移动到曲线上，再水平移动到纵轴，
读出 $f(3) = -1$。}
\end{omfigure}

\begin{method}[读图方法]\label{met:g10:functions:reading}
在 $f$ 的图像上：
\begin{enumerate}
  \item \emph{$a$ 的像}：从横轴上 $x = a$ 处
        竖直走到曲线，再水平走到纵轴，读得的值
        即为 $f(a)$；
  \item \emph{$b$ 的原像}：画出水平直线 $y = b$，
        原像是该直线与曲线全部交点的 $x$
        坐标；
  \item \emph{方程 $f(x) = k$ 的解}：与 $k$ 的原像相同；
  \item \emph{不等式 $f(x) \leq k$ 的解}：所有 $x$，使
        曲线位于直线 $y = k$ 上或其下方。
\end{enumerate}
\end{method}

\begin{omfigure}
\begin{tikzpicture}
\begin{axis}[omaxis,
  width=8.6cm, height=6.0cm,
  xmin=-3.4, xmax=3.4, ymin=-2.4, ymax=4.4,
  xtick={-2,2}, xticklabels={$x_1$,$x_2$}, ytick={2}, yticklabels={$k$}]
  \addplot[omDef, thick, domain=-3:3, samples=120] {0.5*x^2};
  \addplot[omThm, thick, domain=-3.2:3.2] {2};
  \fill (axis cs:-2,2) circle (1.5pt);
  \fill (axis cs:2,2) circle (1.5pt);
  \draw[dashed] (axis cs:-2,2) -- (axis cs:-2,0);
  \draw[dashed] (axis cs:2,2) -- (axis cs:2,0);
\end{axis}
\end{tikzpicture}

{\small 通过图像求解 $f(x) = k$：其解 $x_1$ 与 $x_2$ 是
曲线与水平
直线 $y = k$ 的交点的横坐标。此处 $f(x) \leq k$ 在 $\intcc{x_1}{x_2}$ 上成立，
即曲线位于直线下方。}
\end{omfigure}

\begin{example}\label{ex:g10:functions:membership}
点 $A(2, 5)$ 在 $f(x) = x^2 + 1$ 的图像上吗？计算
$f(2) = 4 + 1 = 5$：是，因为 $f(2)$ 等于 $y$ 坐标，即 $A$ 的纵坐标。
点 $B(3, 8)$ 不在图像上，因为 $f(3) = 10 \neq 8$。
\end{example}

\section{函数的单调性}

\begin{definition}[递增、递减]\label{def:g10:functions:variations}
设 $f$ 在区间 $I$ 上有定义。
\begin{itemize}
  \item 称 $f$ 在 $I$ 上\emph{递增}\index{递增}，当
        对所有 $u, v \in I$，若 $u < v$，则 $f(u) \leq f(v)$：
        输出值随输入值增大而增大，图像从左到右
        上升。
  \item 称 $f$ 在 $I$ 上\emph{递减}\index{递减}，当
        $u < v$ 时有 $f(u) \geq f(v)$：图像从左到右
        下降。
\end{itemize}
若采用严格不等式 $f(u) < f(v)$（相应地 $>$），则称
\emph{严格}递增（相应地递减）。
\end{definition}

\begin{definition}[变化表]\label{def:g10:functions:table}
\emph{变化表}\index{变化表}概括了 $f$ 在哪些
区间上递增与递减，并用箭头 $\nearrow$ 与
$\searrow$ 表示，并记录 $f$ 在转折点处的取值。
\end{definition}

\begin{example}\label{ex:g10:functions:table}
考虑下图所示定义在 $\intcc{-2}{4}$ 上的函数。

\begin{omfigure}
\begin{tikzpicture}
\begin{axis}[omaxis,
  width=8.6cm, height=5.4cm,
  xmin=-2.8, xmax=4.6, ymin=-2.4, ymax=3.8,
  xtick={-2,1,4}, ytick={-1.5,3}, yticklabels={$-1.5$,$3$}]
  \addplot[omDef, thick, domain=-2:4, samples=120] {-0.5*(x-1)^2 + 3};
  \fill (axis cs:-2,-1.5) circle (1.5pt);
  \fill (axis cs:4,-1.5) circle (1.5pt);
  \fill (axis cs:1,3) circle (1.5pt);
\end{axis}
\end{tikzpicture}

{\small 函数定义在 $\intcc{-2}{4}$ 上，先递增后递减。}
\end{omfigure}

由图像可知：$f$ 从 $f(-2) = -1.5$ 递增至其最大值
$f(1) = 3$，然后递减至 $f(4) = -1.5$。其变化表为
\begin{center}
\renewcommand{\arraystretch}{1.3}
\begin{tabular}{c|ccccc}
$x$ & $-2$ & & $1$ & & $4$ \\
\hline
$f$ & $-1.5$ & $\nearrow$ & $3$ & $\searrow$ & $-1.5$ \\
\end{tabular}
\end{center}
\end{example}

\begin{example}[证明单调性]\label{ex:g10:functions:proof}
证明 $f(x) = 3x + 1$ 在 $\R$ 上严格递增。任取
$u < v$，比较它们的像：
\[
f(v) - f(u) = (3v + 1) - (3u + 1) = 3(v - u) > 0,
\]
因为 $v - u > 0$，故 $f(u) < f(v)$，即函数严格
递增。对斜率为负的情形作同样计算，例如\
$g(x) = -2x + 5$，得 $g(v) - g(u) = -2(v-u) < 0$：严格递减。
\end{example}

\section{最值}

\begin{definition}[最大值、最小值]\label{def:g10:functions:extrema}
设 $f$ 定义在 $D$ 上，$a \in D$。函数值 $f(a)$ 称为 $f$ 在 $D$ 上的
\emph{最大值}\index{最大值}，当 $f(x) \leq f(a)$ 对所有
$x \in D$ 成立时；它称为\emph{最小值}\index{最小值}，当
$f(x) \geq f(a)$ 对所有 $x \in D$ 成立时。在图像上，它们分别是曲线的最高点与
最低点。
\end{definition}

\begin{example}\label{ex:g10:functions:extremum}
证明 $f(x) = (x - 1)^2 + 2$ 的最小值为 $2$，定义域为 $\R$，且在
$x = 1$ 处取得。对每个 $x$，平方 $(x-1)^2$ 满足 $\geq 0$，故
$f(x) = (x-1)^2 + 2 \geq 2$；又 $f(1) = 0 + 2 = 2$，因此值 $2$ 确实
能取到。两者合在一起即证明 $2$ 是最小值。
\end{example}

\begin{omfigure}
\begin{tikzpicture}
\begin{axis}[omaxis,
  width=7.6cm, height=5.2cm,
  xmin=-1.8, xmax=3.8, ymin=-0.9, ymax=6.4,
  xtick={1}, ytick={2}]
  \addplot[omDef, thick, domain=-1.0:3.0, samples=100] {(x-1)^2 + 2};
  \addplot[omThm, dashed, domain=-1.6:3.6] {2};
  \fill (axis cs:1,2) circle (1.5pt);
  \node[below right, font=\small] at (axis cs:1,2) {$(1,2)$};
\end{axis}
\end{tikzpicture}

{\small 曲线 $f(x) = (x-1)^2 + 2$ 始终不低于虚线
$y = 2$，并在 $x = 1$ 处与之相切：$f$ 的最小值为 $2$。}
\end{omfigure}

\section{练习}

\begin{exercise}[$\star$]\label{exo:g10:functions:1}
设 $f(x) = 2x^2 - 3x + 1$，求 $0$、$2$、$-1$ 和
$\frac12$ 的像。
\end{exercise}

\begin{exercise}[$\star$]\label{exo:g10:functions:2}
求下列各函数的定义域：
\[
f(x) = 5x - 2, \qquad
g(x) = \frac{3}{x + 4}, \qquad
h(x) = \sqrt{2x - 6}, \qquad
k(x) = \frac{1}{x^2 + 1}.
\]
\end{exercise}

\begin{exercise}[$\star$]\label{exo:g10:functions:3}
设 $f(x) = x^2 - 2x$，求 $0$、$3$ 和 $-1$ 的所有原像。
\end{exercise}

\begin{exercise}[$\star$]\label{exo:g10:functions:4}
点 $A(-1, 4)$ 是否在函数 $f(x) = x^2 - 2x + 1$ 的图像上？
点 $B(2, 1)$ 呢？请通过计算说明理由。
\end{exercise}

\begin{exercise}[$\star$]\label{exo:g10:functions:5}
函数 $g$ 定义在 $\intcc{-3}{5}$ 上，其变化表如下：
\begin{center}
\renewcommand{\arraystretch}{1.3}
\begin{tabular}{c|ccccccc}
$x$ & $-3$ & & $0$ & & $3$ & & $5$ \\
\hline
$g$ & $1$ & $\searrow$ & $-2$ & $\nearrow$ & $4$ & $\searrow$ & $0$ \\
\end{tabular}
\end{center}
\begin{enumerate}
  \item 求 $g$ 在 $\intcc{-3}{5}$ 上的最大值与最小值。
  \item 在不知道具体取值的情况下，比较 $g(-1)$ 与 $g(-0.5)$ 的大小。
  \item 方程 $g(x) = 0$ 在表中每个
        区间上至多有多少个解？
\end{enumerate}
\end{exercise}

\begin{exercise}[$\star\star$]\label{exo:g10:functions:6}
设 $f(x) = -2x + 7$。通过比较 $f(u)$ 与 $f(v)$（$u < v$），
证明 $f$ 在 $\R$ 上严格递减。
\end{exercise}

\begin{exercise}[$\star\star$]\label{exo:g10:functions:7}
证明函数 $f(x) = x^2 + 6x + 11$ 在 $\R$ 上存在最小值，
并求出该最小值及其取得的位置。（提示：写成
$f(x) = (x + 3)^2 + c$ 的形式，求出适当的常数 $c$。）
\end{exercise}

\begin{exercise}[$\star\star$]\label{exo:g10:functions:8}
一个矩形花园的周长为 $40$ m。设其宽为
$x$（单位：m）。
\begin{enumerate}
  \item 分别将长与面积 $A(x)$ 表示为 $x$ 的函数。
        使其有意义的 $x$ 取值范围是什么？
  \item 计算 $A(4)$、$A(8)$、$A(10)$、$A(12)$ 与 $A(16)$，
        并猜想面积最大时矩形的形状。
\end{enumerate}
\end{exercise}

\begin{exercise}[$\star\star$]\label{exo:g10:functions:9}
设 $f(x) = \dfrac{1}{x^2 + 1}$，定义在 $\R$ 上。
\begin{enumerate}
  \item 证明：$0 < f(x) \leq 1$ 对任意实数 $x$ 成立。
  \item $f$ 在 $\R$ 上是否有最大值？是否有最小值？请说明理由。
\end{enumerate}
\end{exercise}

\begin{exercise}[$\star\star\star$]\label{exo:g10:functions:10}
设 $f(x) = x^2$，定义在 $\R$ 上。
\begin{enumerate}
  \item 设 $0 \leq u < v$。对 $f(v) - f(u)$ 作因式分解，由此推出 $f$ 在
        $\intco{0}{+\infty}$ 上严格递增。
  \item 仿照上述论证，证明 $f$ 在
        $\intoc{-\infty}{0}$ 上严格递减。
\end{enumerate}
\end{exercise}

\section{问题：以自身输出为输入的机器}

\begin{problem}[{周末问题 --- 不动点、迭代与
三台著名的机器：海伦的改进器、平方器和
未解的冰雹数}]
\label{pb:g10:functions:1}
函数就是一台机器：输入一个数，输出一个数。
最有趣的实验是把机器的\emph{自身
输出}一次又一次地送回机器，然后两个问题便占据了主导：
这一过程是否会停在某处（\emph{不动点}），以及
机器的变化如何把过程推向那里？本题的机器之一
已为平方根打磨了两千年；
另一台则守着初等数学中最著名的
未解难题。

\medskip
\noindent\textbf{第一部分 --- 海伦的机器。}
考虑函数
\[
f(x) = \frac12\left(x + \frac2x\right).
\]
\begin{enumerate}
  \item 求 $f$ 的定义域，并计算 $f(1)$、$f(2)$ 和
        $f\!\left(\frac32\right)$，结果用精确分数表示。
  \item 求 $\frac32$ 的所有原像：解方程
        $f(x) = \frac32$（先去分母，并运用
        \cref{pb:g10:algebra:1} 中的方法）。
  \item $f$ 的\emph{不动点}\index{不动点}是指被该机器
        保持不变的数：$f(x) = x$。求出 $f$ 的两个
        不动点。哪位老相识守护着
        正的那一个？
  \item 从 $x = 1$ 出发，把机器的输出再次输入机器，
        连续进行三次，并记录精确结果。你重建了
        初中分册无理数周末问题中的哪一个数列，
        以及用函数语言表述的“海伦配方”是什么？
  \item 在同一张图上，画出 $f$ 的图像（$x > 0$），
        以及直线 $y = x$。不动点在这张图上
        在哪里？（可与初中分册中“两支温度计”
        周末问题作比较：两种刻度下读数相同的温度，
        体现的正是同一个想法。）
\end{enumerate}

\medskip
\noindent\textbf{第二部分 --- 为什么这台机器不会失手。}
\begin{enumerate}[resume]
  \item 设 $0 < u < v$，通过通分证明
        下式成立：
        \[
        f(v) - f(u) = \frac{(v - u)(uv - 2)}{2uv} ,
        \]
        由此推出 $f$ 在
        $\intoo{0}{+\infty}$ 上的单调性：先递减至某点，然后递增 ---
        哪一点是分界点？（所用方法与
        \cref{exo:g10:functions:10}相同。）
  \item 由此可知，在 $\intoo{0}{+\infty}$ 上函数 $f$ 有
        最小值，恰好为 $\sqrt2$，并在
        $\sqrt2$ 处取得。这对机器的\emph{每一个}输出
        （正输入）说明了什么？
  \item 证明：若 $x > \sqrt2$，则
        $\sqrt2 < f(x) < x$：猜得过大时总会得到改进，
        而不会冲过头。（为证 $f(x) < x$，比较 $\frac2x$ 与
        $x$。）
  \item 综合第 6--8 题的结论，画出 $f$ 在
        $\intoo{0}{+\infty}$ 上的变化表，
        （见\cref{def:g10:functions:table}），包括最小值。
  \item 第二台机器：$g(x) = x^2$。它的不动点是
        $0$ 与 $1$（请验证）。对 $g$ 从 $x = 0.9$ 出发迭代四次，
        再从 $x = 1.1$ 出发迭代四次（用计算器，
        保留三位小数）。其中一个不动点吸引、另一个排斥：
        哪一个吸引、哪一个排斥？
\end{enumerate}

\medskip
\noindent\textbf{第三部分 --- 冰雹机。}
在正整数上定义：$h(n) = \frac n2$（$n$ 为偶数），
而 $h(n) = 3n + 1$（$n$ 为奇数）。数在 $h$ 的作用下跳动，
像暴风云中的冰雹一样。
\begin{enumerate}[resume]
  \item 计算 $7$ 在 $h$ 作用下的完整“飞行”，直到
        $1$ 为止。需要多少步？最高飞到多高？
  \item 从 $27$ 开始飞行，计算十步。（它的完整飞行需要 $111$ 步，
        最高达到 $9\,232$ --- 起点为 $27$。）
        $7$ 与 $27$ 这两个例子，关于简单规则
        能给我们什么启示？
  \item 解释为什么任何到达 $2$ 的幂的飞行
        都会直接坠落到 $1$，并说明在 $1$ 处发生什么
        （计算 $h(1)$、$h(4)$、$h(2)$）。$h$ 在正整数中
        有没有不动点？每条已观察到的飞行，其末尾
        反而形成的是什么？
  \item \emph{角谷猜想}断言：每一个起始数
        最终都会到达 $1$。计算机已经验证
        了远超 $10^{20}$ 的情形，但尚无人证明它。哪两类
        发现能够\emph{推翻}它？为什么
        大量验证仍然不能了结这一问题
        （初中分册中的“火柴神谕”周末问题已敲响过这一警钟）？
  \item “数学尚未成熟到能回答这类问题，”
        保罗·厄尔多斯如此评论此猜想。请用一句话说明：冰雹机与海伦的
        机器有何区别？在海伦的机器上，
        第 6--8 题已经解决了一切？
\end{enumerate}

\medskip
\noindent\textbf{第四部分 --- 机器使用说明书。}
\begin{enumerate}[resume]
  \item 用区间表示定义域（沿用\cref{pb:g10:numbers:1}中的
        说法）。求下列函数的定义域：
        $x \mapsto \sqrt{x - 3}$, \
        $x \mapsto \frac{1}{x^2 - 9}$, \
        $x \mapsto \sqrt{9 - x^2}$.
  \item 回到海伦的 $f$，求 $3$ 的精确原像
        （通过配方解方程 $f(x) = 3$）。
  \item 一块石子竖直向上抛出，它的高度
        $H(t) = 20t - 5t^2$（单位：米，$t$ 为经过的秒数）。该公式在哪一段时间区间内
        符合实际的物理意义？
        用配方法求最大高度及其对应的
        时刻（见\cref{def:g10:functions:extrema}），并给出
        变化表。
  \item 石子高度至少为 $15$ 米的时间有多长？（用符号表
        解不等式 $H(t) \geq 15$，见
        \cref{met:g10:algebra:signtable}。）
  \item 结语：本题依次考察了四台机器 --- 海伦的改进器、
        平方器、冰雹机和石子的高度。用两三句话说明每个
        函数都具有什么（定义域、法则、图像），以及本题对
        每个函数提出的两个重大问题（迭代最终趋于何处？
        函数如何变化？） --- 这些问题
        将在后续各章中
        逐一讨论。
\end{enumerate}
\end{problem}
